\documentclass{subfiles}
\begin{document}
    Should be clear by now that prime numbers are at the core of cryptography,
        since these are hard to find\footnotemark. On the other hand,
        as in the case of RSA, we need to find distinct prime numbers 
        that are not big, yet not easily computable. 
        Meaning a test for primality would be a good tool.
        Luckly such tool exists -- Pollard's \(p - 1\) algorithm --
        though it's applicable to a specific type of numbers.

    Let's consider the details. Pollard's algorithm is based on the concept of 
        \(B\)-smooth numbers; that is, numbers whose prime factors are all smaller than 
        or equal to a given bound \(B\). Precisely, let \(n = pq\), where \(p, q\) are 
        prime, be the integer to factorize; assume we have found a integer \(L\)
        such that 
        \[
             \divides{p - 1}{L} \text{ and } \ndivides{q - 1}{L}.
        \]
        This means that there exists \(i, j\) and \(k, k \ne 0\) such that 
        \[
            L = i(p - 1) \text{ and } L = j(q - 1) + k.
        \]
        By Fermat little theorem we have 
        \[\begin{aligned}
            a^{L} & \equiv 1 \mod{p} \\ 
            a^{L} & \equiv a^{k} \mod{q}.
        \end{aligned}\]
        Then, we have that 
        \[
            \divides{p}{a^{L} - 1} \text{ and } \ndivides{q}{a^{L} - 1},
        \]
        for most values of \(a\). It follows that \(p = \gcd{a^{L} - 1, n}\).

    The interesting part is that, if \(p - 1\) is such that it's the product of small 
        prime numbers, then it divides \(n!\). In other words, what we do is 
        fix \(a \in \Z^{+}\) and compute \(\gcd{a^{n!}, n} = d\); if \(d \ne 1\) or 
        \(d \ne n\) we have found a non-trivial factor, otherwise we try again by changing the value 
        of \(a\).

    Is easily noted that the algorithm has, at least apparently two severe drawbacks: 
        the computation of the factorial and the consequential complexity. 
        Let us say that, how would be clear later in this section, such drawbacks do not exists.
    
    \footnotetext{The statement is true unless \emph{Riemann} hypothesis turns out to be true.}

    Back to the algorithm: we've said that Pollard's algorithm works if either \(p - 1\) or \(q - 1\)
        are product of small primes; the question is how likely this happens?
        It's here that \(B\)-smoothness plays a key role: if \(n\) is \(B\)-smooth,
        then it divides \(B!\); additionally, any prime \(l\) dividing \(n\) has to be less than or equal to \(B\).
\end{document}
